\documentclass[11pt,a4paper]{article}
\usepackage{gtu-el-yazisi}
\begin{document}
\notbilgi{Elektronik I}{BJT -- 1. ders tekrarı}{BJT-ders1-tekrar.pdf}

\section{npn Junction}
\noindent\sayfa{1}%
\begin{minipage}[c]{0.55\linewidth}
\[
  i_C=\underset{\substack{\downarrow\\ \text{saturation}\\ \text{current}}}{I_S}\,e^{V_{BE}/V_T}
\]
\end{minipage}\hfill
\begin{minipage}[c]{0.4\linewidth}
\centering
\fbox{$\displaystyle I_S=\frac{A_E\,q\,D_n\,n_{p0}}{W}$}

\medskip
\begin{circuitikz}[line width=0.6pt]
  \draw[color=turuncukalem] (0,0) node[npn](Q){};
  \draw[color=turuncukalem] (Q.B) node[left]{$B$};
  \draw[color=turuncukalem] (Q.C) node[right]{$C$};
  \draw[color=turuncukalem] (Q.E) node[left]{$E$} node[right]{$n$};
  \node[text=turuncukalem] at (0,-1.3) {npn};
\end{circuitikz}
\end{minipage}

\[
  i_B=\frac{i_C}{\beta} \qquad \beta \rightarrow \left(\text{influenced by width, and doping}\right)
\]
\[
  i_E=\frac{\beta+1}{\beta}\,i_C=\frac{1}{\alpha}\,i_C=(\beta+1)\,i_B
\]
\[
  i_C=\alpha\,i_E
\]

\sayfa{2}
\begin{center}
\begin{circuitikz}[line width=0.6pt]
  \draw (0,4.6) node[ocirc]{} node[right]{$C$} to[short, i=$i_C$] (0,4) to[cI, l={$I_S\,e^{V_{BE}/V_T}$}] (0,2.4) to[D, l={$D_E\ (I_{SE}=I_S/\alpha)$}] (0,0.6) to[short, i=$i_E$] (0,0) node[ocirc]{};
  \draw (-2.2,2.4) node[ocirc]{} node[left]{$B$} to[short, i=$i_B$] (0,2.4);
  \node at (-0.5,2.0) {$+$}; \node at (-0.75,1.5) {$V_{BE}$}; \node at (-0.5,0.9) {$-$};
\end{circuitikz}
\end{center}

\noindent In npn
\elcizimi[0.4]{sekiller/BJT-ders1-tekrar-s2-1.png}[2]
\begin{align*}
  &\text{Active:} && V_{BE}\geq 0{,}7, && V_{BC}\leq 0{,}4\ \text{V} \\
  &\text{Cut off:} && V_{BE}<0{,}7, && V_{BC}\leq 0{,}4\ \text{V} \\
  &\text{Saturation:} && V_{BE}\geq 0{,}7, && V_{BC}>0{,}4
\end{align*}

\noindent for pnp
\elcizimi[0.35]{sekiller/BJT-ders1-tekrar-s2-2.png}[2]
\begin{align*}
  &\text{Active:} && V_{EB}\geq 0{,}7\ \text{V}, && V_{CB}\leq 0{,}4\ \text{V} \\
  &\text{Cut off:} && V_{EB}<0{,}7\ \text{V}, && V_{CB}\leq 0{,}4\ \text{V} \\
  &\text{Saturation:} && V_{EB}<0{,}7\ \text{V}\ \text{\footnotemark}, && V_{CB}>0{,}4\ \text{V} \quad /\!/
\end{align*}
\footnotetext{Claude notu: muhtemelen doyumda da $V_{EB}\geq 0{,}7$ V olmalı (npn'deki $V_{BE}\geq 0{,}7$ ile tutarlı).}

\sayfa{3}
\noindent
\begin{minipage}[c]{0.3\linewidth}
$2<4$\\
$-2>-4$
\end{minipage}
\begin{minipage}[c]{0.6\linewidth}
$V_C-V_B\leq 0{,}4$ V\\
$V_B-V_C\geq -0{,}4$
\end{minipage}

\section{Operation in Saturation Mode}
\[
  i_C=I_S\,e^{V_{BE}/V_T}-\underbrace{I_{SC}\,e^{V_{BC}/V_T}}_{\substack{V_{BC}>0{,}4\\ \text{için önem taşır}}}
\]
base current in sat region:
\[
  i_B=\frac{I_S}{\beta}\,e^{V_{BE}/V_T}+I_{SC}\,e^{V_{BC}/V_T}
\]
\[
  \text{forced }\beta=\left.\frac{i_C}{i_B}\right|_{sat}\leq\beta
\]
\hrule
\medskip
\noindent $V_{CB}<0{,}4$ olmalı

\begin{ornek}
\textcolor{turuncukalem}{Determine the value of $R_1$ such that $Q_1$ operate in active mode}

\noindent
\begin{minipage}[c]{0.45\linewidth}
\centering
\begin{circuitikz}[line width=0.6pt]
  \draw[color=kirmizikalem] (0,0) node[pnp](Q){};
  \draw[color=kirmizikalem] (Q.E) -- ++(0,0.2) coordinate(EE) to[R, l_=$R_1$] ++(0,1.8) node[above]{$5$ V};
  \draw[color=kirmizikalem] (EE) -- ++(0.9,0) node[circ]{} node[right]{$E$\ \ $0{,}7$ V};
  \draw[color=kirmizikalem] (Q.B) -- ++(-0.6,0) node[ground]{};
  \path (Q.B) ++(-0.6,-0.75) node[text=kirmizikalem]{$B$};
  \draw[color=kirmizikalem] (Q.C) -- ++(0,-0.2) coordinate(CC) to[R, l=$3\,\text{k}\Omega$] ++(0,-1.8) node[below]{$-5$ V};
  \draw[color=kirmizikalem] (CC) -- ++(0.9,0) node[circ]{} node[right]{$C$};
  \node[text=kirmizikalem] at (-1.8,0.3) {$\beta=\infty$};
  \node[text=kirmizikalem] at (-0.8,1) {$0{,}7$ V $+$};
  \node at (1.3,-0.3) {$+0{,}4$};
  \node at (1.4,1.4) {$\downarrow 1{,}8$ mA};
  \node at (1.4,-1.4) {$\downarrow 1{,}8$ mA};
\end{circuitikz}
\end{minipage}\hfill
\begin{minipage}[c]{0.5\linewidth}
\[
  i_E=\frac{5-0{,}7}{R_1}=\frac{4{,}7}{R_1}\ \text{\stepcounter{footnote}\textsuperscript{\thefootnote}}
\]
\[
  i_E=\frac{\beta}{\beta+1}\,i_C=i_C \quad (\beta=\infty)
\]
\[
  V_{CB}=0{,}4\ \text{V}
\]
\[
  i_C=\frac{+0{,}4+5}{3\text{k}}=\frac{5{,}4}{3\text{k}}=1{,}8\ \text{mA}
\]
\end{minipage}
\end{ornek}
\footnotetext[\value{footnote}]{Claude notu: muhtemelen $5-0{,}7=4{,}3$ olmalı, yani $i_E=4{,}3/R_1$.}
\elcizimi[0.35]{sekiller/BJT-ders1-tekrar-s3-1.png}[3]

\sayfa{4}
\[
  R_1\geq\frac{4{,}7}{1{,}8\ \text{mA}}\geq 2{,}3\ \text{k}\Omega\ \text{\footnotemark}
\]
\footnotetext{Claude notu: $4{,}7/1{,}8\approx 2{,}6$; $2{,}3$--$2{,}4$ k$\Omega$ sonucu $4{,}3/1{,}8$ ile tutarlı (bkz. önceki sayfadaki not).}
\noindent
\begin{minipage}[t]{0.3\linewidth}
çünkü
\end{minipage}
\begin{minipage}[t]{0.5\linewidth}
$V_{CB}<0{,}4$\\
eğer akım artarsa\\
$V_{CB}$ artar\\
yani direnç azalırsa
\end{minipage}

\medskip
\noindent Yani direnç artarsa akım de azalır ve öylece çözmüş oluruz. /\!/

\end{document}
